\chapter*{General Conclusion and Perspectives}
\addcontentsline{toc}{chapter}{General Conclusion and Perspectives}
\markboth{General Conclusion and Perspectives}{General Conclusion and Perspectives}

This report addressed a question raised by the modernization of the Energy
Transition Optimizer, an internal BCG~X product for the pre-feasibility analysis
of Power-to-X projects: how can generative AI make an expert-only techno-economic
optimizer accessible in natural language, without compromising the
reproducibility, auditability, and human control its results demand? The answer
defended throughout is a separation between assistance and authority: the
deterministic engine remains the sole source of numerical truth, while a governed
multi-agent layer helps users design, run, and understand a study, with every
state-changing action validated and explicitly approved before it takes effect.

The four scoping questions of Chapter~\ref{chap:context} can now be answered
directly. \emph{What can be delegated:} language models proved well suited to
preparing inputs and explaining outputs, topology drafting, assumption ingestion,
scenario framing, and result interpretation, while every numerical result stayed
with the deterministic engine. \emph{How assisted actions are made trustworthy:}
each state change is a typed, validated proposal that a human approves, and each
assisted turn is traced, so actions are auditable and, taking effect only on
approval, reversible by rejection. \emph{How a non-deterministic capability is
measured:} an independent LLM-as-a-judge scores interpretations on four explicit
dimensions, and a graph-edit-distance metric scores generated structures, both
tracked across versions rather than assumed. \emph{How the work stays adaptable
across clients:} the layer was built as a reusable product capability on a
reproducible, infrastructure-as-code platform, so the same core is adapted per
client rather than rebuilt.

Conducted as a team, the project delivered this on two fronts. It restored the
platform foundation, with federated enterprise identity, reproducible delivery,
and operational observability. And it introduced the governed multi-agent layer
itself, a Director that routes each request to specialist agents for ingestion,
topology design, scenario exploration, and interpretation, behind typed proposals
and human approval. The industrial value of the core this layer serves was
independently confirmed while the work was under way, when the optimizer's
optimization method was granted a United States patent, presented in the
preceding chapter.

At the heart of that layer stand three pillars: a
\textbf{read-only interpretation agent} that explains a simulation population in
plain language while keeping every figure traceable to a deterministic result; a
\textbf{visualization layer}, with interactive scatter exploration and a navigable
discounted-cash-flow table, that makes results legible and carries the
interpretation on demand; and an \textbf{evaluation harness} that measures the
grounding and soundness of interpretations, with a graph-edit-distance metric for
structured outputs. Together they let a non-specialist obtain and audit an
interpretation that previously required an expert. On the small, first-order sample
we evaluated, each assisted stage of a study fell by roughly an order of magnitude
in expert effort, with result interpretation dropping from about two hours to about
twelve minutes (Chapter~\ref{chap:implementation}); these figures are indicative
rather than controlled, and a study with non-specialist participants remains future
work.

The main lesson is that trust in an assisted analytical tool cannot rest on
prompting alone: it depends on the architecture (the read/write split, the approval
gates, the grounding rules), on evidence (measuring the assistant rather than
assuming it), on operations (end-to-end tracing and observability), and ultimately
on the human who approves each state change. Safety, on this view, is a property of
that whole arrangement, and one demonstrated on the analyses we ran rather than
proven in general. Natural extensions follow the same logic: richer comparative
interpretation, broader and human-calibrated evaluation coverage, adversarial
testing of the governance boundaries, and adaptation of these capabilities to each
new client deployment.

On a personal level, the internship strengthened my skills in AI orchestration,
in language-model evaluation and observability, in data visualization, and in
techno-economic financial literacy. Above all, it taught me the discipline it
takes for generative AI to earn trust in a tool whose value depends on being
right, and that is the standard I intend to carry into what I build next.
